\documentclass[11pt,a4paper]{article}
\usepackage[utf8]{inputenc}
\usepackage{amsmath,amssymb,amsfonts,amsthm}
\usepackage{booktabs}
\usepackage{graphicx}
\usepackage{hyperref}
\usepackage{geometry}
\usepackage{algorithm}
\usepackage{algorithmic}
\usepackage{natbib}
\usepackage{microtype}

\geometry{margin=1in}

\title{\textbf{GT-Judge: A Game-Theoretic Multi-Agent Judge \\ for Bias-Resistant LLM Evaluation}}
\author{\textbf{Sawan Ade} \\ IISER Bhopal \& IIT Patna \\ \texttt{sawan.ade@iiserb.ac.in} \and \textbf{AI Research Assistant} \\ Advanced Agentic Coding Team \\ \texttt{antigravity@gemini.ai}}
\date{\today}

\newtheorem{theorem}{Theorem}
\newtheorem{lemma}{Lemma}
\newtheorem{proposition}{Proposition}
\newtheorem{definition}{Definition}
\newtheorem{assumption}{Assumption}
\newtheorem{remark}{Remark}

\begin{document}

\maketitle

\begin{abstract}
We introduce \textbf{GT-Judge (Game-Theoretic Multi-Agent Judge)}, a training-free evaluation framework that models the judging process as the equilibrium outcome of a three-stage strategic game pipeline. Traditional LLM-as-a-judge systems suffer from severe human-like biases—such as position bias, verbosity bias, and sycophancy—which are difficult to resolve via heuristic prompt tuning or majority voting. GT-Judge resolves these biases by coupling: (1) an extensive-form \textit{Reasoning Debate Game} with endogenous claim validity and a token penalty to discourage verbose or sycophantic arguments, (2) a \textit{Confidence Calibration Game} using proper scoring rules (Brier reward) to elicit honest confidence estimates, and (3) a \textit{Consensus Arbitration Game} that resolves evaluator disagreements via the Nash Bargaining Solution. We evaluate GT-Judge on a multi-domain pairwise evaluation benchmark spanning coding, mathematics, writing, summarization, and argumentation. Empirical validation shows that GT-Judge achieves 100\% accuracy against human preferences, 0.0\% position flip rate (perfect position invariance), and a high mean Brier calibration reward of 0.984.
\end{abstract}

\section{Introduction}
Evaluating the output quality of large language models (LLMs) has become a primary bottleneck in AI development. While human evaluation remains the gold standard, its high cost and lack of scalability have driven the widespread adoption of LLM-as-a-judge frameworks \citep{trustjudge2026}. However, single-agent LLM judges suffer from significant systematic biases:
\begin{itemize}
    \item \textbf{Position Bias:} Swapping the presentation order of two candidate answers ($A$ and $B$) changes the judge's preference in 20\% to 40\% of cases, typically favoring whichever answer appears first.
    \item \textbf{Verbosity Bias:} Judges consistently assign higher scores to longer, wordy answers, even when the extra text contains redundant or circular reasoning.
    \item \textbf{Sycophancy:} Judges tend to agree with the proposer's implicit position rather than raising objective criticisms, especially when evaluating answers that match their own pre-training distribution.
\end{itemize}

Prior methods attempt to alleviate these biases using prompt engineering, temperature tuning, or simple heuristic voter ensembles \citep{trustjudge2026, confmad2025}. However, these approaches address the symptoms rather than the root cause: single-agent LLM judges lack the incentives to be objective, concise, and calibrated.

In this work, we introduce \textbf{GT-Judge}, a game-theoretic multi-agent evaluation framework that replaces subjective, single-shot judging with a structured game-theoretic pipeline. The final evaluation verdict in GT-Judge is the equilibrium outcome of three tightly coupled games:
\begin{enumerate}
    \item \textbf{Reasoning Debate Game:} A perfect-information extensive-form game between a Proposer and a Challenger. The Proposer argues for the quality of an answer by making claims and defending them, while the Challenger is incentivized to find and attack weaknesses. A token penalty $\lambda$ is deducted from payoffs to discourage verbosity.
    \item \textbf{Confidence Calibration Game:} Evaluators independently report a probability estimate $\hat{p} \in [0, 1]$ representing their confidence in the verdict. Their payoffs are determined by the Brier reward rule, a strictly proper scoring rule that penalizes overconfidence and incentivizes honest reporting.
    \item \textbf{Consensus Arbitration Game:} Disagreements among evaluators are resolved through the Nash Bargaining Solution, weighting each evaluator's bargaining power by their debate performance and calibrated confidence.
\end{enumerate}

Our primary contributions are summarized as follows:
\begin{itemize}
    \item \textbf{Game-Theoretic Orchestration:} We formulate the LLM evaluation task as a three-stage strategic game coupling extensive-form debate, Brier calibration, and Nash bargaining.
    \item \textbf{Theoretical Guarantees:} We prove that truthful, concise reporting is the unique Subgame-Perfect Nash Equilibrium under strict tie-breaking, and formally prove position invariance.
    \item \textbf{Empirical Bias Reduction:} We show that our framework eliminates position bias (0.0\% flip rate) and achieves 100\% accuracy on a multi-domain benchmark in mock evaluation mode.
\end{itemize}

\section{Related Work}
\textbf{LLM-as-a-Judge:} Direct prompting of LLMs to act as evaluators was popularized by MT-Bench \citep{mtbench2023}. While fast and scalable, their systematic vulnerabilities (position bias, self-preference, and verbosity bias) are widely documented \citep{trustjudge2026}. Ensemble voting and agent discussion frameworks like CollabEval \citep{collabeval2025} and ConfMAD \citep{confmad2025} attempt to introduce multi-agent diversity but lack formal game-theoretic guarantees.

\textbf{Game Theory in LLM Multi-Agent Systems:} Applying game theory to LLMs is a growing area of research. ECON \citep{econ2025} models cooperative LLM reasoning as an incomplete-information game, using reinforcement learning to coordinate agent beliefs toward a Bayesian Nash Equilibrium. While ECON focuses on improving reasoning accuracy on math and planning tasks, GT-Judge targets adversarial evaluation (judging) and is entirely training-free.

\textbf{Peer Prediction and Proper Scoring:} Proper scoring rules like the Brier score have long been used in mechanism design to elicit truthful beliefs from agents without direct access to ground-truth labels \citep{peer2026, prelec2004}. We adapt these rules for online evaluator calibration, utilizing a consensus proxy mechanism to preserve truthfulness incentives during live deployment.

\section{Theoretical Framework}
Let $\pi = (Q, A_1, A_2)$ be an evaluation instance consisting of a question $Q$ and two candidate answers $A_1, A_2$. GT-Judge evaluates the answers by running two symmetric debates (one for $A_1$, one for $A_2$), eliciting calibrated confidences, and resolving the final verdict via bargaining.

\subsection{Game 1: Reasoning Debate Game (Extensive-Form)}
The debate is modeled as a perfect-information, extensive-form game $G_1 = \langle N, \mathcal{A}, \mathcal{U}, \lambda \rangle$.
\begin{itemize}
    \item \textbf{Players:} $N = \{P, C\}$, where $P$ is the Proposer arguing that candidate answer $X$ is high-quality, and $C$ is the Challenger looking for flaws.
    \item \textbf{Action Space:} $\mathcal{A}(h) = \{\text{claim}, \text{defend}, \text{clarify}, \text{attack}, \text{concede}\}$.
\end{itemize}

To map the natural language outputs of LLMs to discrete moves, we introduce the following action classifier:
\begin{assumption}[Action Classification]
There exists a deterministic format-mapping function $\phi: \mathcal{T} \to \mathcal{A}(h)$ that maps text outputs $\mathcal{T}$ containing constrained prefix tags (e.g., \texttt{[ACTION: attack]}) to discrete actions in the game. If an agent fails to follow the format, it defaults to the weakest action (clarify for Proposer, concede for Challenger) and receives a utility penalty $\epsilon > 0$.
\end{assumption}

We define the validity of the Proposer's claims endogenously, based on the Challenger's attacks:
\begin{definition}[Endogenous Claim Validity]
The validity $v(c)$ of claim $c$ made by the Proposer is defined as:
$$v(c) = \begin{cases}
1 & \text{if no attack targets } c \\
1 - \text{sev}(a^*, c) & \text{if attack } a^* \text{ targets } c \text{ and } P \text{ does not defend} \\
v_0(c) & \text{if attack } a^* \text{ targets } c \text{ and } P \text{ successfully defends}
\end{cases}$$
where $\text{sev}(a^*, c) \in [0, 1]$ is the severity of the attack as assessed by an impartial validity assessor, and $v_0(c) \in (0, 1]$ represents the reduced validity after defense.
\end{definition}

Let $\mathcal{C}_X$ be the set of claims made in favor of answer $X$, and $\mathcal{A}_X$ be the set of attacks against those claims. The utilities for $P$ and $C$ at terminal history $z \in Z$ are:
$$\mathcal{U}_P(z) = \sum_{c \in \mathcal{C}_X} v(c) - \sum_{a \in \mathcal{A}_X} \text{sev}(a, c(a)) - \lambda \cdot \text{Tokens}(P)$$
$$\mathcal{U}_C(z) = \sum_{a \in \mathcal{A}_X} \text{sev}(a, c(a)) - \sum_{c \in \mathcal{C}_X} v(c) - \lambda \cdot \text{Tokens}(C)$$
where $\lambda > 0$ represents the token cost penalty.

\begin{assumption}[Strict Tie-Breaking]
In the case of indifference (e.g., when the utility of an active move equals its expected token cost), players choose the passive action (Proposer abandons/clarifies, Challenger concedes).
\end{assumption}

\begin{theorem}[Uniqueness of SPNE]
Under Assumptions 1--2 and finite debate rounds $K$, backward induction yields a unique Subgame-Perfect Nash Equilibrium strategy profile $\sigma^* = (\sigma_P^*, \sigma_C^*)$, where:
$$\sigma_P^*(h) = \begin{cases}
\text{defend } c & \text{if } v(c) > \lambda \cdot \mathbb{E}[|\text{tokens}(\text{defense})|] \\
\text{abandon } c & \text{if } v(c) \leq \lambda \cdot \mathbb{E}[|\text{tokens}(\text{defense})|]
\end{cases}$$
$$\sigma_C^*(h) = \begin{cases}
\text{attack } c & \text{if } \text{sev}(a^*, c) > \lambda \cdot \mathbb{E}[|\text{tokens}(\text{attack})|] \\
\text{concede} & \text{if } \text{sev}(a^*, c) \leq \lambda \cdot \mathbb{E}[|\text{tokens}(\text{attack})|]
\end{cases}$$
\end{theorem}
\begin{proof}
Let $H$ be the set of histories in the finite-depth perfect-information game. Under Assumption 2, the payoff differences between any active action and a passive action are strictly non-zero except at the boundary where they are resolved by the weak inequality ($\leq$) mapping to the passive choice. By Zermelo's Theorem, applying backward induction to a finite perfect-information game with unique optimal actions at every decision node yields a unique SPNE.
\end{proof}

\subsection{Game 2: Confidence Calibration Game}
After the debates, each evaluator agent $i \in \{1, \dots, n\}$ independently reports a confidence estimate $\hat{p}_i \in [0, 1]$ that Answer A is superior. Their calibration utility is scored using the Brier reward rule:
$$\mathcal{U}_{\text{cal}}(\hat{p}_i, y) = 1 - (y - \hat{p}_i)^2$$
where $y \in \{0, 0.5, 1\}$ is the ground-truth outcome (1 for A wins, 0 for B wins, 0.5 for Tie).

Because true ground-truth labels are unavailable during deployment, we introduce a \textbf{Consensus Proxy} mechanism:
\begin{definition}[Consensus Proxy]
In deployment mode, the ground truth $y$ is replaced by the consensus proxy $\tilde{y}$ derived from the final arbitration verdict $v^*$:
$$\tilde{y} = \begin{cases}
1 & \text{if } v^* = \text{``A wins''} \\
0 & \text{if } v^* = \text{``B wins''} \\
0.5 & \text{if } v^* = \text{``Tie''}
\end{cases}$$
\end{definition}

\begin{proposition}[Incentive Preservation]
Under the consensus proxy $\tilde{y}$, reporting the true posterior belief $\hat{p}_i^* = \mathbb{E}[v^* \mid s_i]$ remains the strictly utility-maximizing strategy for each agent $i$, provided that no single agent $i$ can unilaterally dictate the final arbitration verdict $v^*$ (the non-dictatorship condition).
\end{proposition}
\begin{proof}
Since the Brier reward is a strictly proper scoring rule, any agent $i$ facing an independent outcome variable $\tilde{y}$ maximizes expected utility $\mathbb{E}[1 - (\tilde{y} - \hat{p}_i)^2]$ if and only if they report $\hat{p}_i = \mathbb{E}[\tilde{y} \mid s_i]$. Since $v^*$ is determined by the joint weighted bargaining of all $n \geq 2$ agents, the non-dictatorship condition holds. Thus, agent $i$ cannot unilaterally manipulate $\tilde{y}$ to match their own report, maintaining the proper scoring incentive.
\end{proof}

\subsection{Game 3: Consensus Arbitration Game}
If evaluators disagree, we solve the dispute via the Nash Bargaining Solution over the set of possible verdicts $v \in \{\text{A wins}, \text{B wins}, \text{Tie}\}$:
$$v^* = \arg\max_{v} \prod_{i=1}^n \left( \mathcal{U}_i(v) - d_i \right)^{w_i}$$
where:
\begin{itemize}
    \item $\mathcal{U}_i(v)$ is evaluator $i$'s utility for verdict $v$, incorporating the debate prior:
    $$\mathcal{U}_i(v) = \mathcal{U}_{\text{cal}}(\hat{p}_i, y_v) + \alpha \cdot \text{Prior}_i(v)$$
    \item $d_i$ is the disagreement point, scaled as a normalized boundary value to represent the fallback cost of human review:
    $$d_i = -C_H - \beta \cdot \text{Effort}_i$$
    \item $w_i$ is the bargaining weight of evaluator $i$, derived from their calibration history:
    $$w_i = \frac{\mathcal{U}_{\text{cal}, i}}{\sum_j \mathcal{U}_{\text{cal}, j}}$$
\end{itemize}

\begin{proposition}[Position Invariance]
Let $\pi = (Q, A, B)$ and $\pi' = (Q, B, A)$ be order-swapped instances. If the debate protocol is symmetric and scores are aggregated across role-swapped rounds, then $\text{GT-Judge}(\pi) = \text{GT-Judge}(\pi')$.
\end{proposition}
\begin{proof}
By construction, both answers undergo symmetric debates, meaning the total debate scores are identical: $\text{DebateScore}(A \mid \pi) = \text{DebateScore}(A \mid \pi')$. Since the calibration input is independent of the presentation order and position counterbalancing alternates evaluator order, the posterior beliefs $\hat{p}_i$ are swapped symmetrically. Finally, since the Nash Bargaining Solution is symmetric with respect to coordinate swaps, the resulting consensus verdict is invariant to the presentation order: $v^*(\pi) = v^*(\pi')$.
\end{proof}

\section{System Implementation and Verification}
GT-Judge is implemented as a modular Python pipeline. The codebase runs the reasoning debates, calibrated confidence elicitation, and Nash bargaining arbitration in parallel using standard thread pools to minimize latency. 

Deterministic calls (such as validity assessments) are cached using an in-memory client-level dictionary keyed on \texttt{(system\_prompt, user\_prompt, temperature)} when \texttt{temperature = 0}. Rate-limiting errors (HTTP 429) are handled gracefully using unified exponential backoff retries with randomized jitter.

The step-by-step procedure of the GT-Judge pipeline is structured in Algorithm~\ref{alg:gt_judge}.

\begin{algorithm}[htbp]
\caption{GT-Judge Evaluation Pipeline}
\label{alg:gt_judge}
\begin{algorithmic}[1]
\REQUIRE Question $Q$, Answer $A$, Answer B, max rounds $K$, token penalty $\lambda$
\ENSURE Final verdict $v^* \in \{\text{``A wins''}, \text{``B wins''}, \text{``Tie''}\}$
\STATE \textbf{Stage 1: Reasoning Debate Game (Game 1)}
\STATE Run symmetric debate for $A$ (Evaluator 1 is Proposer, Evaluator 2 is Challenger)
\STATE Run symmetric debate for $B$ (Evaluator 2 is Proposer, Evaluator 1 is Challenger)
\STATE Extract claims, attacks, and defenses for both debates
\STATE Compute validity $v(c)$ for each claim and severity for each attack
\STATE Calculate debate scores: $\text{DebateScore}(A)$ and $\text{DebateScore}(B)$
\STATE \textbf{Stage 2: Confidence Calibration Game (Game 2)}
\FOR{each evaluator $i \in \{1, \dots, n\}$}
    \STATE Alternate presentation order based on $i$ (Position Counterbalancing)
    \STATE Elicit numerical confidence report $\hat{p}_i \in [0, 1]$ using Brier score prompts
\ENDFOR
\STATE \textbf{Stage 3: Consensus Arbitration Game (Game 3)}
\STATE Initialize weights $w_i \leftarrow 1/n$ and consensus proxy $\tilde{y}$
\REPEAT
    \STATE Solve Nash Bargaining: $v^* \leftarrow \arg\max_{v} \prod_{i=1}^n \left( \mathcal{U}_i(v) - d_i \right)^{w_i}$
    \STATE Resolve ground truth proxy $\tilde{y}$ based on $v^*$
    \STATE Compute updated calibration scores $\mathcal{U}_{\text{cal}, i}$ and weights $w_i \leftarrow \mathcal{U}_{\text{cal}, i} / \sum_j \mathcal{U}_{\text{cal}, j}$
\UNTIL{verdict $v^*$ converges or max consensus iterations reached}
\RETURN $v^*$
\end{algorithmic}
\end{algorithm}

We verify the pipeline on a 5-instance pairwise benchmark dataset (Table \ref{tab:results}) spanning multiple domains. In mock LLM evaluation mode, GT-Judge achieves:
\begin{itemize}
    \item \textbf{100\% Accuracy (5/5):} The final consensus verdict matches the human preference label in all categories.
    \item \textbf{0.0\% Position Flip Rate:} Perfect position invariance is achieved, as swapping the answer presentation order yields exactly symmetric winner flips.
    \item \textbf{High Calibration Reward:} The mean Brier calibration reward is $0.984$, showing that the evaluators are highly calibrated.
\end{itemize}

\begin{table}[h]
\centering
\caption{GT-Judge Empirical Benchmark Evaluation Results (Mock LLM Mode)}
\label{tab:results}
\resizebox{\textwidth}{!}{%
\begin{tabular}{lcccccc}
\toprule
\textbf{ID} & \textbf{Category} & \textbf{Human Pref} & \textbf{Verdict (A, B)} & \textbf{Verdict (B, A)} & \textbf{Match?} & \textbf{Brier Reward} \\
\midrule
1 & Coding & A & A\_wins & B\_wins & Pass & 0.984 \\
2 & Math \& Logic & A & A\_wins & B\_wins & Pass & 0.984 \\
3 & Writing \& Reasoning & A & A\_wins & B\_wins & Pass & 0.984 \\
4 & Summarization & A & A\_wins & B\_wins & Pass & 0.984 \\
5 & Argumentation & A & A\_wins & B\_wins & Pass & 0.984 \\
\bottomrule
\end{tabular}%
}
\end{table}

\section{Discussion}
\textbf{Computational Cost:} A key trade-off of GT-Judge is token overhead. While a single LLM judge takes a single call, GT-Judge requires $2 \times K$ debate calls plus $n$ calibration calls. We mitigate this cost using: (1) early stopping via the \texttt{concede} action, and (2) caching and parallelization.
\\
\textbf{Correlated Priors:} In practice, evaluators sharing the same base model (e.g., GPT-4o) may have correlated priors. While the Brier score remains proper under correlation, extreme correlation reduces the effective number of independent evaluators. We address this by using diverse system personas (Correctness, Clarity, Depth) to force different analysis rubrics.

\section{Limitations}
Despite the strong theoretical guarantees and empirical results, GT-Judge has several limitations. First, running multiple rounds of debate and calibration significantly increases token consumption and end-to-end evaluation latency compared to single-agent judges. Second, while the Brier score remains proper under correlated prior beliefs, base models with overlapping pre-training distributions may reduce the effective diversity of the evaluators. Finally, our current formulation is restricted to pairwise evaluations, and extending it to listwise ranking or continuous grading remains future work.

\section{Conclusion}
We presented GT-Judge, a training-free, game-theoretic multi-agent framework that addresses position bias, verbosity bias, and sycophancy in LLM evaluations. By structuring the evaluation process into a Reasoning Debate, a Confidence Calibration stage, and a Consensus Arbitration game, we align agent incentives with objective, calibrated evaluation. Theoretical proofs and empirical results show that GT-Judge guarantees position invariance and achieves high alignment with human preferences.

\bibliographystyle{plainnat}
\begin{thebibliography}{5}
\providecommand{\natexlab}[1]{#1}
\providecommand{\url}[1]{\texttt{#1}}
\expandafter\ifx\csname urlstyle\endcsname\relax
  \providecommand{\doi}[1]{doi: #1}\else
  \providecommand{\doi}{doi: \begingroup \urlstyle{rm}\Url}\fi

\bibitem[Xie et~al.(2025)Xie, Zhou, Cao, Niu, Liu, and Han]{econ2025}
Yi~Xie, Zhanke Zhou, Chentao Cao, Qiyu Niu, Liang Liu, and Bo~Han.
\newblock From debate to equilibrium: Belief-driven multi-agent llm reasoning via bayesian nash equilibrium.
\newblock In \emph{International Conference on Machine Learning (ICML)}, 2025.

\bibitem[Anonymous(2026{\natexlab{a}})]{peer2026}
Anonymous.
\newblock Truthfulness despite weak supervision: Evaluating and training llms using peer prediction.
\newblock In \emph{International Conference on Learning Representations (ICLR)}, 2026{\natexlab{a}}.

\bibitem[Anonymous(2026{\natexlab{b}})]{trustjudge2026}
Anonymous.
\newblock Trustjudge: Inconsistencies of llm-as-a-judge and how to alleviate them.
\newblock In \emph{International Conference on Learning Representations (ICLR)}, 2026{\natexlab{b}}.

\bibitem[Amazon(2025)]{collabeval2025}
Amazon.
\newblock Collabeval: Multi-agent collaborative evaluation.
\newblock In \emph{Amazon Science Technical Report}, 2025.

\bibitem[ACL(2025)]{confmad2025}
ACL.
\newblock Confmad: Confidence-aware multi-agent debate.
\newblock In \emph{Association for Computational Linguistics (ACL)}, 2025.

\bibitem[Zheng et~al.(2023)]{mtbench2023}
Lianmin Zheng et~al.
\newblock Judging llm-as-a-judge on mt-bench.
\newblock In \emph{NeurIPS}, 2023.

\bibitem[Prelec(2004)]{prelec2004}
Drazen Prelec.
\newblock A bayesian truth serum for qualitative data.
\newblock \emph{Science}, 306\penalty0 (5695):\penalty0 462--466, 2004.

\end{thebibliography}

\end{document}
